
%%%%%%%%%%%%%%%%%%%%%%%%%%%%%%%%%%%%%%%%
%           main body
%%%%%%%%%%%%%%%%%%%%%%%%%%%%%%%%%%%%%%%%


%-----------------------------------
%	TITLE PAGE
%-----------------------------------

\title[高等代数]{\LARGE \bfseries 高等代数} % The short title appears at the bottom of every slide, the full title is only on the title page
\author[伪逆与最小二乘法] % Your institution as it will appear on the bottom of every slide, may be shorthand to save space
{\Large \bfseries 伪逆与最小二乘法}
% \author{Singular Value Decomposition} % Your name
% \date{\today} % Date, can be changed to a custom date
\date{}

%------------------------------------------------

\begin{document}


%------------------------------------------------

\begin{frame}

\titlepage % Print the title page as the first slide

\end{frame}

%------------------------------------------------

\begin{frame}
\frametitle{问题的引入}

对于解线性方程组$A\mathbf{x} = \mathbf{b}$的问题, 我们知道当$A$是可逆方阵的时候, 该线性方程组有唯一解$\mathbf{x} = A^{-1}\mathbf{b}$.

\vspace{1em}

当$A$没有逆 (是不可逆方阵, 或者不是方阵) 的时候, 有什么补救办法呢?

\pause
\vspace{1em}

希望: 定义一个伪逆 (Pseudoinverse), 记作$A^+$, 使得当$A$可逆的时候, $A$的伪逆就是它的逆, 即
\[A^+ = A^{-1}\]
如果可能的话, 还要求当$A\mathbf{x} = \mathbf{b}$无解时, $\mathbf{x} = A^+\mathbf{b}$是``近似程度''最好的``解''.

\end{frame}

%------------------------------------------------

% \begin{frame}
% \frametitle{内容提要} % Table of contents slide, comment this block out to remove it
% \tableofcontents % Throughout your presentation, if you choose to use \section{} and \subsection{} commands, these will automatically be printed on this slide as an overview of your presentation
% \end{frame}

%------------------------------------------------

\section{广义逆}

%------------------------------------------------

\subsection{广义逆的定义与性质}

%------------------------------------------------

\begin{frame}
\frametitle{广义逆的定义与性质}

\begin{Def}[矩阵的广义逆, Generalized Inverse]
$m\times n$ 矩阵 $A$ 的广义逆指的是一个 $n\times m$ 的矩阵 $X,$ 满足
\[AXA = A.\]
通常把$A$的广义逆记作$A^-.$
\end{Def}

\pause

\begin{remark}[广义逆的普遍存在性]
任意矩阵 $A$ 的广义逆 $A^-$ 总存在. 设 $P,Q$ 分别为 $m$ 阶与 $n$ 阶可逆阵, 使得 $A = P \begin{pmatrix} I_r & 0 \\ 0 & 0 \end{pmatrix}_{m\times n}\!\!\!\!\!\!\!\!\!\!\! Q^{-1},$ 那么
\[A^- = Q \begin{pmatrix} I_r & * \\ *' & *'' \end{pmatrix}_{n\times m}\!\!\!\!\!\!\!\!\!\!\! P^{-1}\]
其中 $r = r(A)$, $*,*',*''$ 分别是大小为 $r \times (m-r)$, $(n-r) \times r,$ 以及 $(n-r) \times (m-r)$ 的块, 块中元素可以任取.
\end{remark}

\end{frame}

%------------------------------------------------

\begin{frame}
\frametitle{广义逆的定义与性质}

可以通过 $A, A^-$ 的分块矩阵表达式看出:
\begin{itemize}
\item $AA^- = P \begin{pmatrix} I_r & * \\ 0 & 0 \end{pmatrix} P^{-1}.$
\item $A^-A = Q \begin{pmatrix} I_r & 0 \\ *' & 0 \end{pmatrix} Q^{-1}.$
\item $AA^-A = P \begin{pmatrix} I_r & 0 \\ 0 & 0 \end{pmatrix} Q^{-1} = A.$
\item $A^-AA^- = Q \begin{pmatrix} I_r & * \\ *' & *' \cdot * \end{pmatrix} P^{-1}.$
\end{itemize}

\end{frame}

%------------------------------------------------

\begin{frame}
\frametitle{广义逆的几何解释}

记 $A = \begin{pmatrix} \gamma_1 \\ \vdots \\ \gamma_m \end{pmatrix} = \begin{pmatrix} \alpha_1 & \cdots & \alpha_n \end{pmatrix}$, 其中 $\gamma_i \in \mathbb{R}^n$ 是 $A$ 的行向量, $\alpha_j \in \mathbb{R}^m$ 是 $A$ 的列向量. $A$ 可以视作一个从 $\mathbb{R}^n$ 到 $\mathbb{R}^m$ 的线性变换
\[
\sigma: \mathbb{R}^n \to \mathbb{R}^m, \quad \mathbf{x} \mapsto A\mathbf{x}.
\]
那么有空间的正交分解
\[
\mathbb{R}^n = C(A^T) \Otop N(A), \quad \mathbb{R}^m = C(A) \Otop N(A^T),
\]
其中
\begin{itemize}
\item $N(A) = \{\mathbf{x}\in \mathbb{R}^n: A\mathbf{x} = \mathbf{0}\} = \{\mathbf{x}\in \mathbb{R}^n: \langle \mathbf{x}, \gamma_i \rangle = 0, i=1,\ldots,m\}$ 是 $A$ 的 (右) 零空间, 也是 $\sigma$ 的核 $\ker \sigma;$
\item $C(A^T) = \operatorname{span}\{\gamma_1, \ldots, \gamma_m\}$ 是 $A$ 的行空间;
\item $N(A^T) = \{\mathbf{y}\in \mathbb{R}^m: A^T\mathbf{y} = \mathbf{0}\} = \{\mathbf{y}\in \mathbb{R}^m: \langle \mathbf{y}, \alpha_j \rangle = 0, j=1,\ldots,n\}$ 是 $A$ 的左零空间;
\item $C(A) = \operatorname{span}\{\alpha_1, \ldots, \alpha_n\}$ 是 $A$ 的列空间, 也是 $\sigma$ 的值域 $\operatorname{im} \sigma.$
\end{itemize}

\pause

\begin{remark}
若 $\mathbb{R}^n, \mathbb{R}^m$ 分别以 $Q, P$ 的列向量为基, 则 $\sigma$ 的矩阵表示为 $\begin{pmatrix} I_r & 0 \\ 0 & 0 \end{pmatrix}.$ 这时, $N(A)$ 和 $C(A)$ 分别由 $Q$ 的后 $n-r$ 列, $P$ 的前 $r$ 列张成.
\end{remark}

\end{frame}

%------------------------------------------------

\begin{frame}
\frametitle{广义逆的几何解释}

\begin{block}{一些记号}
约定几个投影映射:
\begin{gather*}
\operatorname{pr}_1: \mathbb{R}^n = C(A^T) \Otop N(A) \to C(A^T), \quad \mathbf{x} = \mathbf{x}_1 + \mathbf{x}_2 \mapsto \mathbf{x}_1;\\
\operatorname{pr}_2: \mathbb{R}^n = C(A^T) \Otop N(A) \to N(A), \quad \mathbf{x} = \mathbf{x}_1 + \mathbf{x}_2 \mapsto \mathbf{x}_2;\\
\operatorname{pr}_1: \mathbb{R}^m = C(A) \Otop N(A^T) \to C(A), \quad \mathbf{y} = \mathbf{y}_1 + \mathbf{y}_2 \mapsto \mathbf{y}_1;\\
\operatorname{pr}_2: \mathbb{R}^m = C(A) \Otop N(A^T) \to N(A^T), \quad \mathbf{y} = \mathbf{y}_1 + \mathbf{y}_2 \mapsto \mathbf{y}_2.
\end{gather*}
\end{block}
约定记号 $\operatorname{id} = \operatorname{pr}_1 \Odown \operatorname{pr}_2,$ 进而任何一个线性映射 $f: \mathbb{R}^k \to \mathbb{R}^m$ 都可以写成如下的形式:
\[
f = f_1 \Odown f_2 = (\operatorname{pr}_1 \circ f) \Odown (\operatorname{pr}_2 \circ f): \mathbb{R}^k \to \mathbb{R}^m \xrightarrow{\operatorname{pr}_1 \Odown \operatorname{pr}_2} C(A) \Otop N(A^T).
\]
\alert{注意, 这里并不是说 $f$ 的值域 $f(\mathbb{R}^k)$ 包含 $C(A)$ 或者 $N(A^T)$ 的子空间, 而是说 $f$ 的值域中的每一个元素都可以唯一分解成 $C(A)$ 中的一个元素与 $N(A^T)$ 中的一个元素的和.}
我们采用类似的记号, 记 $g_1: W \to \mathbb{R}^m,$ $g_2: W^\perp \to \mathbb{R}^m$ 得到的线性映射为
\[
g_1 \Otop g_2: W \Otop W^\perp \to \mathbb{R}^m, \quad \mathbf{w} + \mathbf{w}' \mapsto g_1(\mathbf{w}) + g_2(\mathbf{w}').
\]

\end{frame}

%------------------------------------------------

\begin{frame}
\frametitle{广义逆的几何解释}

那么 $\sigma$ 自然地诱导了一个空间的同构
\[
\tilde{\sigma}: C(A^T) \cong \mathbb{R}^n / \ker \sigma \to \operatorname{im} \sigma = C(A), \quad \mathbf{x} + N(A) \mapsto A\mathbf{x}.
\]
由于 $\tilde{\sigma}$ 是一个同构, 它的逆 $\tilde{\sigma}^{-1}: C(A) \to C(A^T)$ 也是一个同构. 于是我们可以定义 $A$ 的广义逆 $A^-$ 使得
\[
\tau: \mathbb{R}^m \to \mathbb{R}^n, \quad \mathbf{y} \mapsto A^-\mathbf{y}
\]
满足
\begin{itemize}
\item $\tau|_{C(A)} = \tilde{\sigma}^{-1} \Odown \eta$, 其中 $\eta: C(A) \to N(A)$ 是任意一个线性映射;
\item 对左零空间 $N(A^T)$ 中的任意 $\mathbf{y}$ 映射过去的像没有任何要求, $A^-\mathbf{y}$ 可以任取 (保证是线性映射即可).
\end{itemize}
综合起来, 就是
\[
\tau = (\tilde{\sigma}^{-1} ~ {\color{red}\Odown} ~ \eta) ~ {\color{cyan}\Otop} ~ \zeta : C(A) ~ {\color{cyan}\Otop} ~ N(A^T) \to \mathbb{R}^n = C(A^T) ~ {\color{red}\Otop} ~ N(A),
\]
其中 $\eta: C(A) \to N(A)$ 是任意一个线性映射, $\zeta: N(A^T) \to \mathbb{R}^n$ 是任意一个线性映射. 这样构造, 才能使下面的复合映射还是 $\sigma:$
\[
\mathbb{R}^n \xrightarrow{\sigma} \mathbb{R}^m \xrightarrow{\tau} \mathbb{R}^n \xrightarrow{\sigma} \mathbb{R}^m.
\]


\end{frame}

%------------------------------------------------

\subsection{广义逆与线性方程组的解}

%------------------------------------------------

\begin{frame}
\frametitle{广义逆与齐次线性方程组的解}

\begin{prop}[广义逆与齐次线性方程组]
齐次线性方程组 $A\mathbf{x} = \mathbf{0}$ 的解的全体为
\[\mathbf{x} = (I_n - A^-A)\mathbf{z},\]
其中 $A^-$ 取遍 $A$ 的所有广义逆, $\mathbf{z}$ 取遍所有的 $n$ 维列向量.
\end{prop}

\pause
\vspace{1em}

\startproof[bg=yellow, fg=red]
{\color{red}\ding{43}} 首先证明 $\mathbf{x} = (I_n - A^-A)\mathbf{z}$ 是 $A\mathbf{x} = \mathbf{0}$ 的解: 直接验证
\[
A\mathbf{x} = A(I_n - A^-A)\mathbf{z} = (A - \underbrace{AA^-A}_{=A})\mathbf{z} = \mathbf{0}.
\]

\pause
\vspace{0.5em}

{\color{red}\ding{43}} 反过来, 任意一个 $A\mathbf{x} = \mathbf{0}$ 的解 $\mathbf{x}_0$ 都可以写成 $(I_n - A^-A)\mathbf{z}$ 的形式: 取 $\mathbf{z} = \mathbf{x}_0$, 则
\[
(I_n - A^-A)\mathbf{z} = (I_n - A^-A)\mathbf{x}_0 = \mathbf{x}_0 - A^- \underbrace{A\mathbf{x}_0}_{=0} = \mathbf{x}_0.
\]
\qed

\end{frame}

%------------------------------------------------

\begin{frame}
\frametitle{广义逆与齐次线性方程组的解 | 几何解释}

齐次线性方程组 $A\mathbf{x} = \mathbf{0}$ 的解的全体就是 $N(A) = \ker \sigma$. 之前我们已经知道, $\tau = (\tilde{\sigma}^{-1} \Odown \eta) \Otop \zeta,$ 其中 $\eta: C(A) \to N(A)$ 是任意一个线性映射, $\zeta: N(A^T) \to \mathbb{R}^n$ 是任意一个线性映射. 考虑复合映射
\[
\xymatrix@R=0.3em{
C(A^T) \Otop N(A)
  \ar[r]^-{\sigma}
&
C(A) \Otop N(A^T)
  \ar[r]^-{\tau}
&
C(A^T) \Otop N(A)
\\
\mathbf{z}_1+\mathbf{z}_2
  \ar@{|->}[r]
&
\sigma(\mathbf{z}_1)+\mathbf{0}
  \ar@{|->}[r]
&
(\tilde{\sigma}^{-1}\Odown\eta)(\sigma(\mathbf{z}_1))
+\zeta(\mathbf{0})
}
\]
可知 $(\tau \circ \sigma)(\mathbf{z}_1 + \mathbf{z}_2) = (\tilde{\sigma}^{-1}\Odown\eta)(\sigma(\mathbf{z}_1)) = \mathbf{z}_1 + \eta(\sigma(\mathbf{z}_1))$. 于是
\[
(I_n - A^-A)\mathbf{z} = (\operatorname{id} - \tau \circ \sigma)(\mathbf{z}_1 + \mathbf{z}_2) = \mathbf{z}_2 - \eta(\sigma(\mathbf{z}_1)) \in N(A),
\]
并且可以跑遍 $N(A)$ 中的所有元素 (例如直接固定 $\mathbf{z}_1 = \mathbf{0}$, 让 $\mathbf{z}_2$ 跑遍 $N(A)$ 就行了).

\end{frame}

%------------------------------------------------

\begin{frame}
\frametitle{广义逆与非齐次线性方程组的解}

\begin{prop}[广义逆与非齐次线性方程组]
非齐次线性方程组 $A\mathbf{x} = \mathbf{b}$ 有解, 当且仅当存在 $A$ 的广义逆 $A^-$ 使得
\[\mathbf{b} = A A^- \mathbf{b}.\]
有解时, 它的解的全体为 $A^-\mathbf{b}$, $A^-$ 取遍 $A$ 的所有广义逆.
\end{prop}

\pause
\vspace{0.5em}

代数上的证明与齐次的情况类似, 这里我们给出几何上的解释.

\vspace{0.5em}

还是基于之前的分解 $\tau = (\tilde{\sigma}^{-1} \Odown \eta) \Otop \zeta$, 其中 $\eta: C(A) \to N(A)$ 是任意一个线性映射, $\zeta: N(A^T) \to \mathbb{R}^n$ 是任意一个线性映射. 记 $\mathbf{b} = \mathbf{b}_1 + \mathbf{b}_2,$ 其中 $\mathbf{b}_1 \in C(A), \mathbf{b}_2 \in N(A^T)$. 那么有
\[
AA^-\mathbf{b} = \sigma \circ \tau(\mathbf{b}_1 + \mathbf{b}_2) = \sigma \circ \tau(\mathbf{b}_1) + \sigma \circ \tau(\mathbf{b}_2) = \sigma \circ (\tilde{\sigma}^{-1} \Odown \eta)(\mathbf{b}_1) + \sigma \circ \zeta(\mathbf{b}_2) = \mathbf{b}_1 + \underbrace{\sigma(\zeta(\mathbf{b}_2))}_{\in C(A)}.
\]
所以有 ~~ $A\mathbf{x} = \mathbf{b}$ 有解 ~ $\Longleftrightarrow$ ~ $\mathbf{b}_2 = 0$ ~ $\Longleftrightarrow$ ~ $AA^-\mathbf{b} = \mathbf{b}_1 = \mathbf{b}$.

\pause
\vspace{0.5em}

另一方面, 对于固定的 $\mathbf{b} \in C(A),$ 有
\[
A^- \mathbf{b} = \tau(\mathbf{b}) = (\tilde{\sigma}^{-1} \Odown \eta)(\mathbf{b}) = \tilde{\sigma}^{-1}(\mathbf{b}) + \eta(\mathbf{b}),
\]
当 $\eta$ 跑遍所有 $C(A) \to N(A)$ 的线性映射时, $\eta(\mathbf{b})$ 就跑遍了 $N(A)$, 从而 $A^-\mathbf{b}$ 就跑遍了 $\tilde{\sigma}^{-1}(\mathbf{b}) + N(A)$, 也就是 $A\mathbf{x} = \mathbf{b}$ 的所有解.

\end{frame}

%------------------------------------------------

\begin{frame}
\frametitle{广义逆与线性方程组的解}

{
\larger[1]
\begin{question}
如何用矩阵分块的方法 (代数上的方法) 来证明与解释上面两个命题呢?
\end{question}
}

\end{frame}

%------------------------------------------------

\begin{frame}

{
\larger[1]
广义逆只解决了我们的第一部分要求, 即它``像是''逆, 而且当矩阵的确可逆的时候它就是矩阵通常意义下的逆. 但是它并不唯一, 而且仍然没有解决求``最佳近似解''的需求.
}

\end{frame}

%------------------------------------------------

\section{Moore-Penrose伪逆}

%------------------------------------------------

\subsection{伪逆的定义与性质}

%------------------------------------------------

\begin{frame}
\frametitle{伪逆的定义}

\begin{Def}[Moore-Penrose伪逆, Pseudoinverse]
$m\times n$ 矩阵 $A$ 的 Moore-Penrose 伪逆, 或简称伪逆, 指的是同时满足如下四个条件的 $n\times m$ 矩阵 $X$
\[
\begin{cases}
AXA = A \\
XAX = X \\
(AX)^T = AX \\
(XA)^T = XA
\end{cases}
\]
通常, $A$ 的伪逆被记作 $A^+$.
\end{Def}

\pause
\vspace{0.5em}

\begin{remark}
新增加的后 3 个条件, 用于消去广义逆 $A^- = Q \begin{pmatrix} I_r & * \\ *' & *'' \end{pmatrix} P^{-1}$ 中的三块不确定性, 从而使得 $A^+$ 是唯一的.
\end{remark}
\pause
从代数上看 (分块矩阵的形式), 这个结论是比较容易看出来的 (前提是知道 $AX, XA$ 都是对称幂等阵, 这是马上要证明的). 接下来重点从几何的角度来解释一下这四个条件的意义.

\end{frame}

%------------------------------------------------

\begin{frame}
\frametitle{伪逆的几何解释}

{\color{red}\ding{43}} 由条件 1 $AXA = A$ 或条件 2 $XAX = X$ 可知, $AX$ 以及 $XA$ 都是幂等阵, 从而 $\tau\circ\sigma$ 与 $\sigma\circ\tau$ 都是幂等变换.

\pause
\vspace{0.5em}
对于一个一般的幂等变换 $f: \mathbb{R}^k \to \mathbb{R}^k,$ 由于 $\ker f \cap \operatorname{im} f = \{\mathbf{0}\}$, $f$ 诱导了一个空间的直和分解
\[
\mathbb{R}^k = \operatorname{im} f \oplus \ker f.
\]

\pause
\vspace{1em}

{\color{red}\ding{43}} 由条件 3 $(AX)^T = AX$ 知 $AX$ 是一个对称阵, 从而 $\tau\circ\sigma$ 是 $\mathbb{R}^n$ 上的一个对称幂等变换.
\pause
类似地, 由条件 4 $(XA)^T = XA$ 知 $XA$ 是一个对称阵, 从而 $\sigma\circ\tau$ 是 $\mathbb{R}^m$ 上的一个对称幂等变换.

\pause
\vspace{0.5em}

对于一个一般的对称幂等变换 $f: \mathbb{R}^k \to \mathbb{R}^k,$ $\mathbb{R}^k = \operatorname{im} f \oplus \ker f,$ 任取 $\alpha \in \ker f, \beta = f(\gamma) \in \operatorname{im} f,$ 都有
\[
\langle \alpha, \beta \rangle = \langle \alpha, f(\gamma) \rangle = \langle f(\alpha), \gamma \rangle = \langle \mathbf{0}, \gamma \rangle = 0.
\]
于是 $\operatorname{im} f \perp \ker f$, 即有空间的正交分解 $\mathbb{R}^k = \operatorname{im} f \Otop \ker f.$

\pause
\vspace{0.5em}

于是, 有空间分解
\[
\mathbb{R}^n = \operatorname{im}(\tau\circ\sigma) \Otop \ker(\tau\circ\sigma), \quad \mathbb{R}^m = \operatorname{im}(\sigma\circ\tau) \Otop \ker(\sigma\circ\tau).
\]

\end{frame}

%------------------------------------------------

\begin{frame}
\frametitle{伪逆的几何解释}

\[
\mathbb{R}^n = \operatorname{im}(\tau\circ\sigma) \Otop \ker(\tau\circ\sigma), \quad \mathbb{R}^m = \operatorname{im}(\sigma\circ\tau) \Otop \ker(\sigma\circ\tau).
\]

{\color{red}\ding{43}} 考察 $\tau\circ\sigma: \mathbb{R}^n \to \mathbb{R}^n$.

\pause
\vspace{0.5em}

任取 $\mathbf{x} = \mathbf{x}_1 + \mathbf{x}_2 \in \mathbb{R}^n,$ 其中 $\mathbf{x}_1 \in N(A), \mathbf{x}_2 \in C(A^T)$.

由 $\mathbf{x}_1 \in N(A)$ 知 $(\tau\circ\sigma)(\mathbf{x}_1) = \tau(\mathbf{0}) = \mathbf{0}$, 从而 $N(A) \subseteq \ker(\tau\circ\sigma)$.

\pause
\vspace{0.5em}

反之, 若 $(\tau\circ\sigma)(\mathbf{x}) = \mathbf{0}$, 即 $XA\mathbf{x} = \mathbf{0}$, 由条件 1 ($AXA = A$) 知
\[
A\mathbf{x} = AXA\mathbf{x} = A\underbrace{(XA\mathbf{x})}_{=\,\mathbf{0}} = \mathbf{0},
\]
故 $\mathbf{x} \in N(A)$. 从而 $\ker(\tau\circ\sigma) = N(A)$.

\pause
\vspace{0.5em}

再由 $\tau\circ\sigma$ 的对称性 (条件 4), $\operatorname{im}(\tau\circ\sigma) = \ker(\tau\circ\sigma)^\perp = N(A)^\perp = C(A^T)$. 于是, $\tau\circ\sigma$ 就是 $\mathbb{R}^n$ 上到行空间 $C(A^T)$ 的\textbf{正交投影}:
\[
\tau\circ\sigma = \operatorname{pr}_1: \mathbb{R}^n = C(A^T) \Otop N(A) \to C(A^T), \quad (\tau\circ\sigma)(\mathbf{x}_1 + \mathbf{x}_2) = \mathbf{0} + \mathbf{x}_2 = \mathbf{x}_2.
\]

\end{frame}

%------------------------------------------------

\begin{frame}
\frametitle{伪逆的几何解释}

{\color{red}\ding{43}} 类似地考察 $\sigma\circ\tau: \mathbb{R}^m \to \mathbb{R}^m$. 对任意 $A\mathbf{x} \in C(A)$, 由条件 1:
\[
(\sigma\circ\tau)(A\mathbf{x}) = AX(A\mathbf{x}) = (AXA)\mathbf{x} = A\mathbf{x},
\]
从而 $C(A) \subseteq \operatorname{im}(\sigma\circ\tau)$; 又显然 $\operatorname{im}(\sigma\circ\tau) \subseteq C(A)$, 故 $\operatorname{im}(\sigma\circ\tau) = C(A)$. 再由 $\sigma\circ\tau$ 的对称性 (条件 3):
\[
\ker(\sigma\circ\tau) = \operatorname{im}(\sigma\circ\tau)^\perp = C(A)^\perp = N(A^T).
\]
于是, $\sigma\circ\tau$ 是 $\mathbb{R}^m$ 上到列空间 $C(A)$ 的\textbf{正交投影}:
\[
\sigma\circ\tau = \operatorname{pr}_1: \mathbb{R}^m = C(A) \Otop N(A^T) \to C(A), \quad (\sigma\circ\tau)(\mathbf{y}_1 + \mathbf{y}_2) = \mathbf{y}_1 + \mathbf{0} = \mathbf{y}_1.
\]

\end{frame}

%------------------------------------------------

\begin{frame}
\frametitle{伪逆的几何解释}

{\color{red}\ding{43}} 由广义逆 $\tau = (\tilde{\sigma}^{-1} \Odown \eta) \Otop \zeta$ 知 (取 $\mathbf{z}_1 \in C(A^T), \mathbf{z}_2 \in N(A)$):
\[
(\tau\circ\sigma)(\mathbf{z}_1 + \mathbf{z}_2) = \mathbf{z}_1 + \eta(\sigma(\mathbf{z}_1)).
\]
而 $\tau\circ\sigma = \operatorname{pr}_1$ 给出 $(\tau\circ\sigma)(\mathbf{z}_1 + \mathbf{z}_2) = \mathbf{z}_1$, 故
\[
\eta(\sigma(\mathbf{z}_1)) = \mathbf{0}, \quad \forall\, \mathbf{z}_1 \in C(A^T).
\]
由于 $\tilde{\sigma} = \sigma|_{C(A^T)}: C(A^T) \overset{\sim}{\to} C(A)$ 是同构 (特别地是满射), 故 $\eta = \mathbf{0}$.

\pause
\vspace{0.5em}

{\color{red}\ding{43}} 有了 $\eta = \mathbf{0}$, 对任意 $\mathbf{b}_2 \in N(A^T)$ 有 $\tau(\mathbf{b}_2) = \zeta(\mathbf{b}_2)$. 由 $\ker(\sigma\circ\tau) = N(A^T)$:
\[
\sigma(\zeta(\mathbf{b}_2)) = (\sigma\circ\tau)(\mathbf{b}_2) = \mathbf{0} \implies \zeta(\mathbf{b}_2) \in N(A).
\]
另一方面, 由条件 2 ($XAX = X$, 即 $\tau\circ\sigma\circ\tau = \tau$):
\[
\zeta(\mathbf{b}_2) = \tau(\mathbf{b}_2) = (\tau\circ\sigma)\left(\tau(\mathbf{b}_2)\right) = \operatorname{pr}_1\left(\zeta(\mathbf{b}_2)\right) \implies \zeta(\mathbf{b}_2) \in C(A^T).
\]
综合 $\zeta(\mathbf{b}_2) \in N(A)$ 与 $\zeta(\mathbf{b}_2) \in C(A^T)$, 由 $C(A^T) \cap N(A) = \{\mathbf{0}\}$ 得 $\zeta = \mathbf{0}$.

\pause
\vspace{0.5em}

因此, 伪逆 $A^+$ 对应的 $\tau$ 由下式唯一确定:
\[
\tau = (\tilde{\sigma}^{-1} ~ {\color{red}\Odown} ~ 0) ~ {\color{cyan}\Otop} ~ 0 : C(A) ~ {\color{cyan}\Otop} ~ N(A^T) \to \mathbb{R}^n = C(A^T) ~ {\color{red}\Otop} ~ N(A),
\]
即 $A^+$ 将列空间 $C(A)$ 通过 $\tilde{\sigma}^{-1}$ 同构地映回行空间 $C(A^T)$, 而将左零空间 $N(A^T)$ 映射到 $\mathbf{0}$.

\end{frame}

%------------------------------------------------

\begin{frame}
\frametitle{伪逆的性质}

\begin{prop}[伪逆的存在性与唯一性]
任意 $m\times n$ 矩阵 $A$ 的伪逆 $A^+$ 都存在, 而且
\[A^+ = V\Sigma^+U^T,\]
其中 $A = U\Sigma V^T$ 为矩阵 $A$ 的奇异值分解, 奇异值为 $\sigma_1, \ldots, \sigma_r$,
\[\Sigma^+ = \begin{pmatrix}
\frac{1}{\sigma_1} & & & \\ & \ddots & & \\ & & \frac{1}{\sigma_r} & \\ & & & \bigzero \end{pmatrix}_{n\times m}\]
\end{prop}

\pause
\vspace{1em}

\startproof[bg=yellow, fg=red]
直接验证 $X = V\Sigma^+U^T$ 满足四个条件即可. 例如, 验证第一个条件:
\[
AXA = (U\Sigma V^T)(V\Sigma^+U^T)(U\Sigma V^T) = U\Sigma\Sigma^+\Sigma V^T = U\Sigma V^T = A.
\]

\end{frame}

%------------------------------------------------

\begin{frame}
\frametitle{伪逆的性质}

\begin{remark}[伪逆的另一种表达]
$A$ 的奇异值分解又可以写作\(\displaystyle A = U\Sigma V^T = \sum\limits_{i=1}^r \sigma_i \mathbf{u}_i \mathbf{v}_i^T \). 于是我们又有 $A^+$ 的另一种表达
\[A^+ = \sum\limits_{i=1}^r \sigma_i^{-1} \mathbf{v}_i \mathbf{u}_i^T\]
其中这些 $\mathbf{v}_i, \mathbf{u}_i$ 分别是正交阵 $V$ 与 $U$ 的列.
\end{remark}

\end{frame}

%------------------------------------------------

\subsection{伪逆与线性方程组}

%------------------------------------------------

\begin{frame}
\frametitle{伪逆与线性方程组}

\begin{prop}[伪逆与线性方程组解的关系]
对于线性方程组 $A\mathbf{x} = \mathbf{b}$ ($\mathbf{b}$ 可以是零向量也可以不是), 它若有解, 则解可以由系数矩阵 $A$ 的 Moore-Penrose 伪逆统一给出:
\[A^+\mathbf{b} + (I - A^+A)\mathbf{w}\]
其中 $\mathbf{w}$ 是任意一个 $n$ 维列向量. 当然, $A\mathbf{x} = \mathbf{b}$ 有解的充要条件为
\[AA^+\mathbf{b} = \mathbf{b}\]
\end{prop}
这里是类似于``特解+通解''的形式, 其中 $A^+\mathbf{b}$ 是一个特解, $(I - A^+A)\mathbf{w}$ 是通解.

\pause
\vspace{1em}

证明这里就不写了, 与广义逆的情况是类似的.

\end{frame}

%------------------------------------------------

\section{最小二乘法}

%------------------------------------------------

\begin{frame}
\frametitle{最小二乘法: 问题分析}

我们来考虑非齐次线性方程组 $A\mathbf{x} = \mathbf{b}$ 无解的情况, $A$ 是一个 $m\times n$ 矩阵. 此时, 我们有
\[A\mathbf{x} = \mathbf{b}\text{无解} \Longleftrightarrow \mathbf{b}\not\in C(A)\]

转而: 在 $C(A)$ 中寻找 $A\widehat{\mathbf{x}}$, 使得它与 $\mathbf{b}$ 最接近, 即求
\[\argmin_{\widehat{\mathbf{x}}}\lVert A\widehat{\mathbf{x}} - \mathbf{b} \rVert,\]
并称其为一个最小二乘解.

\vspace{2em}

可以将以上结果与之前的广义逆相关的结果对比来看.

\end{frame}

%------------------------------------------------

\begin{frame}
\frametitle{最小二乘法: 法方程组}

{\bfseries 事实}: $\mathbb{R}^m$ 中的一个点, 与一个线性子空间中距离最短的点是这个点在这个子空间上的投影 (见下图).

\begin{figure}[H]
\centering
\begin{tikzpicture}
\path[-, thin] (0,0) edge (7,0);
\path[-, thin] (0,0) edge (-2,-2);
\path[-, thin] (-2,-2) edge (5,-2);
\path[-, thin] (7,0) edge (5,-2);
\path[-Stealth, very thick] (0,0) edge (4,2);
\path[-Stealth, very thick, dashed] (4,2) edge (4,-1);
\path[-Stealth, very thick] (0,0) edge (4,-1);
\path[-Stealth, very thick] (0,0) edge (2.5,-1.5);
\path[-Stealth, very thick, dashed] (4,2) edge (2.5,-1.5);
\path[-Stealth, dotted] (4,-1) edge (2.5,-1.5);
\node at (4.7, 0.5) {$\beta-\alpha$};
\node at (2.6, 0.5) {$\beta'-\alpha$};
\node at (1, 0.7) {$\alpha$};
\node at (2.3, -0.28) {$\beta$};
\node at (0.6, -0.7) {$\beta'$};
\node at (-1, -1.6) {$W$};
\node at (4.1, -1.4) {$\beta'-\beta$};
\draw[-, very thick] (4,-0.6) -- (3.6, -0.5) -- (3.6, -0.9);
\end{tikzpicture}
\end{figure}

\vspace{1em}
\pause

{\bfseries Hint}: 考察 $\beta'-\alpha, \beta-\alpha, \beta'-\beta$ 围成的直角三角形的边长关系.

\end{frame}

%------------------------------------------------

\begin{frame}
\frametitle{最小二乘法: 法方程组}

于是
\begin{align*}
\lVert A\widehat{\mathbf{x}} - \mathbf{b} \rVert ~ \text{取得最小值} ~ & \Longleftrightarrow A\widehat{\mathbf{x}} - \mathbf{b} \perp C(A) \\
& \Longleftrightarrow A^T (A\widehat{\mathbf{x}} - \mathbf{b}) = \mathbf{0} \\
& \Longleftrightarrow A^TA\widehat{\mathbf{x}} = A^T\mathbf{b}
\end{align*}

那么, 问题转化为了求解非齐次线性方程组
\[A^TA\widehat{\mathbf{x}} = A^T\mathbf{b}.\]

以上方程组 (被称作法方程组, Normal Equations) 的特点:
\begin{itemize}
\item 总有解；
\item 即使解 $\widehat{\mathbf{x}}$ 不唯一, 投影 $A\widehat{\mathbf{x}}$ 是唯一的.
\end{itemize}

\end{frame}

%------------------------------------------------

\begin{frame}
\frametitle{最小二乘法}

\begin{block}{最小二乘解}
\begin{itemize}
\item<1-> 若 $r(A) = n$, 那么 $r(A^T) = r(A) = n$. 此时 $A^TA$ 可逆, $A^TA\widehat{\mathbf{x}} = A^T\mathbf{b}$ 有唯一解. 于是我们得到唯一的最小二乘解
\[\widehat{\mathbf{x}} = (A^TA)^{-1}A^T\mathbf{b}.\]
\item<2-> 若 $r(A) < n$, 那么 $r(A^T) = r(A) < n$. 此时 $A^TA$ 不可逆, $A^TA\widehat{\mathbf{x}} = A^T\mathbf{b}$ 有无穷多组解, 亦即最小二乘解有无穷多组.
\end{itemize}
\end{block}

\pause
\vspace{1em}

\begin{question}
在 $A^TA$ 不可逆的情况下, 如何从无穷多组最小二乘解中选出一个``最佳''的解?
\end{question}

\end{frame}

%------------------------------------------------

\begin{frame}
\frametitle{最小二乘解不唯一的几何直观}

$r(A) < n$ 时, $N(A) \neq \{\mathbf{0}\}$, 映射 $\mathbf{x}\mapsto A\mathbf{x}$ 不是单射:
不同的 $\mathbf{x}$ 可以给出相同的 $A\mathbf{x}$.

\vspace{0.3em}

\begin{columns}[T]
\begin{column}{0.52\textwidth}
\begin{itemize}
  \item $\mathbf{b}$ 在 $C(A)$ 上的正交投影 $\hat{\mathbf{b}}$ 仍然\textbf{唯一}
  \item 满足 $A\hat{\mathbf{x}} = \hat{\mathbf{b}}$ 的 $\hat{\mathbf{x}}$ 不唯一: 若 $\mathbf{z}\in N(A)$,
  \[A(\hat{\mathbf{x}} + \mathbf{z}) = A\hat{\mathbf{x}} + \underbrace{A\mathbf{z}}_{=\mathbf{0}} = \hat{\mathbf{b}}\]
  \item 全体最小二乘解构成仿射子空间
  \[\hat{\mathbf{x}}_0 + N(A)\]
\end{itemize}
\vspace{0.3em}
在其中取\textbf{范数最小}者, 即完全落在行空间 $C(A^T)$ 中的那个, 就是 $\mathbf{x}^+ = A^+\mathbf{b}$.\\[0.2em]
{\small (图: 蓝线方向 $\parallel N(A)$, 虚线方向 $\parallel C(A^T)$, 两者正交, 即 $C(A^T)\perp N(A)$)}
\end{column}
\begin{column}{0.44\textwidth}
\centering
\scalebox{1.25}{
\begin{tikzpicture}[scale=1.5]
  \draw[->] (-0.2, 0) -- (2.6, 0) node[right, scale=0.75] {$x_1$};
  \draw[->] (0, -0.2) -- (0, 2.6) node[above, scale=0.75] {$x_2$};
  % \foreach \x in {1,2} {
  %   \draw (\x, 0.05) -- (\x,-0.05) node[below, scale=0.65] {$\x$};
  %   \draw (0.05, \x) -- (-0.05,\x) node[left, scale=0.65] {$\x$};
  % }
  % Solution set (affine subspace) -- direction (1,-1), i.e. N(A)
  \draw[blue, thick] (-0.1, 2.1) -- (2.1, -0.1);
  \node[blue, scale=0.65] at (2.2, 0.35) {$\hat{\mathbf{x}}_0\!+\!N(A)$};
  % N(A) direction arrow
  \draw[-Stealth, blue!80, thin] (0.6,1.4) -- (0.85,1.15)
      node[right, scale=0.6, yshift=16pt, xshift=-10pt] {$N(A)$};
  % Other solutions
  \fill[gray!60] (0.3, 1.7) circle (1.5pt);
  \fill[gray!60] (1.8, 0.2) circle (1.5pt);
  % Perpendicular from origin to min-norm solution -- direction (1,1), i.e. C(A^T)
  \draw[dashed, red!80!black] (0,0) -- (1,1);
  \node[red!80!black, scale=0.6, rotate=45] at (0.38, 0.52) {$C(A^T)$};
  % Right angle mark at (1,1): sides along (1,1)/sqrt2 and (1,-1)/sqrt2
  % epsilon/sqrt2 ≈ 0.085 with epsilon=0.12
  \draw[red!80!black, thin]
      (0.92, 0.92) -- (1.00, 0.83) -- (1.08, 0.92);
  % Origin
  \fill[red!80!black] (0,0) circle (2pt) node[below left, scale=0.7] {$\mathbf{0}$};
  % Min norm solution
  \fill[red!80!black] (1,1) circle (2.5pt)
      node[above right, xshift=1.5pt, scale=0.7] {$\mathbf{x}^+$};
\end{tikzpicture}
}
\end{column}
\end{columns}

\end{frame}

%------------------------------------------------

\begin{frame}
\frametitle{最小二乘解与伪逆的关系}

\begin{thm}[最小二乘解与伪逆的关系]
$\mathbf{x}^+ := A^+\mathbf{b}$ 是长度最小的最小二乘解. 其中 $A^+$ 为矩阵 $A$ 的 Moore-Penrose 伪逆.
\end{thm}

\vspace{1em}
\pause

\startproof[bg=yellow, fg=red]
以下设 $A$ 的奇异值分解为 $A = U\Sigma V^T.$

\vspace{1em}

{\color{red}\ding{43}} $\mathbf{x}^+$ 是最小二乘解:

\pause
\vspace{0.5em}

将其代入法方程组 $A^TA\widehat{\mathbf{x}} = A^T\mathbf{b}$ 左边有
\begin{align*}
A^TA\cdot\mathbf{x}^+ & = A^TAA^+ \cdot \mathbf{b} \\
& = (V\Sigma^TU^T)(U\Sigma V^T)(V\Sigma^+ U^T) \mathbf{b} = V\Sigma^T\Sigma\Sigma^+ U^T \cdot \mathbf{b} \\
& = V \begin{pmatrix} \Sigma_0^T & 0 \\ 0 & 0 \end{pmatrix} \begin{pmatrix} \Sigma_0 & 0 \\ 0 & 0 \end{pmatrix} \begin{pmatrix} \Sigma_0^{-1} & 0 \\ 0 & 0 \end{pmatrix} U^T \cdot \mathbf{b} \\
& = V \Sigma^T U^T \cdot \mathbf{b} = A^T \cdot \mathbf{b} = \text{法方程组右边}
\end{align*}
其中 $\Sigma_0 = diag(\sigma_1, \ldots, \sigma_r).$ 由上式即知, $\mathbf{x}^+$ 是最小二乘解.

\end{frame}

%------------------------------------------------

\begin{frame}
\frametitle{最小二乘解与伪逆的关系}

{\color{red}\ding{43}} $\mathbf{x}^+$ 是最小二乘解中长度最小者:

\pause
\vspace{0.5em}

设 $\widehat{\mathbf{x}}_0$ 也是一个最小二乘解, 即满足 $A^TA\widehat{\mathbf{x}}_0 = A^T\mathbf{b}.$ 于是
\begin{align*}
& A^TA(\widehat{\mathbf{x}}_0-\mathbf{x}^+) = A^T\mathbf{b} - A^T\mathbf{b} = \mathbf{0} \\
\Longrightarrow & \lVert A (\widehat{\mathbf{x}}_0-\mathbf{x}^+) \rVert^2 = (\widehat{\mathbf{x}}_0-\mathbf{x}^+)^T A^TA (\widehat{\mathbf{x}}_0-\mathbf{x}^+) = 0 \\
\Longrightarrow & A (\widehat{\mathbf{x}}_0-\mathbf{x}^+) = \mathbf{0}
\end{align*}

\pause

$A$ 用其奇异值分解又可写作 $A = \sum\limits_{i=1}^r \sigma_i \mathbf{u}_i \mathbf{v}_i^T$. 于是从上式我们有
\begin{align*}
A (\widehat{\mathbf{x}}_0-\mathbf{x}^+) = \mathbf{0} & \Longrightarrow \sum\limits_{i=1}^r \sigma_i \mathbf{u}_i \underbrace{\color{red} \mathbf{v}_i^T (\widehat{\mathbf{x}}_0-\mathbf{x}^+)}_{\textbf{这是个数}} = \mathbf{0} \\
& \Longrightarrow \sum\limits_{i=1}^r \left( \sigma_i \cdot \left( \mathbf{v}_i^T (\widehat{\mathbf{x}}_0-\mathbf{x}^+) \right) \right) \mathbf{u}_i = \mathbf{0}
\end{align*}

\end{frame}

%------------------------------------------------

\begin{frame}
\frametitle{最小二乘解与伪逆的关系}

但是, 我们知道 $\mathbf{u}_1,\ldots,\mathbf{u}_r$ 是线性无关的 (甚至是相互正交的), 而且 $\sigma_1,\ldots,\sigma_r$ 都是大于0的实数. 于是
\[\mathbf{v}_1^T (\widehat{\mathbf{x}}_0-\mathbf{x}^+) = \cdots = \mathbf{v}_r^T (\widehat{\mathbf{x}}_0-\mathbf{x}^+) = 0.\]
那么
\begin{align*}
\langle \mathbf{x}^+, \widehat{\mathbf{x}}_0 - \mathbf{x}^+ \rangle & = (\mathbf{x}^+)^T (\widehat{\mathbf{x}}_0 - \mathbf{x}^+) = (A^+\mathbf{b})^T (\widehat{\mathbf{x}}_0 - \mathbf{x}^+) \\
& = \mathbf{b}^T (A^+)^T (\widehat{\mathbf{x}}_0 - \mathbf{x}^+) \\
& = \mathbf{b}^T \left( \sum\limits_{i=1}^r \sigma_i^{-1} \mathbf{v}_i \mathbf{u}_i^T \right)^T (\widehat{\mathbf{x}}_0 - \mathbf{x}^+) \\
& = \mathbf{b}^T \left( \sum\limits_{i=1}^r \sigma_i^{-1} \mathbf{u}_i \underbrace{\mathbf{v}_i^T (\widehat{\mathbf{x}}_0 - \mathbf{x}^+)}_{=\mathbf{0}} \right) = 0
\end{align*}

因此, 我们有
\[\lVert \widehat{\mathbf{x}}_0 \rVert^2 = \lVert \mathbf{x}^+ + (\widehat{\mathbf{x}}_0 - \mathbf{x}^+) \rVert^2 = \lVert \mathbf{x}^+ \rVert^2 + \lVert (\widehat{\mathbf{x}}_0 - \mathbf{x}^+) \rVert^2 \geqslant \lVert \mathbf{x}^+ \rVert^2\]
也就是说 $\mathbf{x}^+ = A^+\mathbf{b}$ 的确是长度最小的最小二乘解.
\qed

\end{frame}

%------------------------------------------------

\begin{frame}
\frametitle{最小二乘解与伪逆的关系}

\begin{remark}
{\color{red}注意}! 最小二乘解 $\mathbf{x}^+ = A^+\mathbf{b}$ 包含了 $A^TA$ 可逆时的唯一最小二乘解 $(A^TA)^{-1}A^T\mathbf{b}$.
\end{remark}

\vspace{1em}
\pause

\startproof[bg=yellow, fg=red]
根据矩阵乘积的秩关系
\[A^TA\text{可逆} \Longrightarrow n = r(A^TA) = r(A) = r(\Sigma),\]
于是, $\Sigma$ 就必须形如 $\Sigma = \begin{pmatrix} \Sigma_0 \\ 0 \end{pmatrix}$, 其中 $\Sigma_0 = diag(\sigma_1, \ldots, \sigma_n)$. 于是
\[A^+ = V\begin{pmatrix} \Sigma_0^{-1} & 0 \end{pmatrix} U^T\]
另一方面,
\begin{align*}
(A^TA)^{-1}A^T & = (V\Sigma^TU^TU\Sigma V^T)^{-1} V\Sigma^TU^T \\
& = V \Sigma_0^{-2} V^T V\Sigma^TU^T = V\begin{pmatrix} \Sigma_0^{-1} & 0 \end{pmatrix} U^T
\end{align*}
\qed

\end{frame}

%------------------------------------------------

\begin{frame}
\frametitle{例: 最小二乘解不唯一}

设 $A = \begin{pmatrix}1 & 1 & 0 \\ 1 & 1 & 0 \\ 0 & 0 & 1\end{pmatrix}$, $\mathbf{b} = \begin{pmatrix}0\\4\\2\end{pmatrix}$.
注意 $r(A) = 2 < n = 3$, $N(A) = \mathrm{span}\{(1,-1,0)^T\}$, $\mathbf{b}\notin C(A)$.

\pause
\vspace{0.5em}

\textbf{法方程组} $A^TA\hat{\mathbf{x}} = A^T\mathbf{b}$:
\[A^TA = \begin{pmatrix}2&2&0\\2&2&0\\0&0&1\end{pmatrix}, \quad A^T\mathbf{b} = \begin{pmatrix}4\\4\\2\end{pmatrix}\]

化简后得 $x_1+x_2 = 2$, $x_3 = 2$.
{\color{red}3 个未知量, 2 个独立约束 $\to$ 1 个自由变量, 无穷多解!}

\pause
\vspace{0.3em}

全体最小二乘解: $\left\{(t,\; 2-t,\; 2)^T : t\in\mathbb{R}\right\}$.

\end{frame}

%------------------------------------------------

\begin{frame}
\frametitle{例: 用伪逆求最小范数最小二乘解}

$A$ 的 Moore-Penrose 伪逆为
\[A^+ = \begin{pmatrix}\dfrac{1}{4} & \dfrac{1}{4} & 0 \\[6pt] \dfrac{1}{4} & \dfrac{1}{4} & 0 \\[2pt] 0 & 0 & 1\end{pmatrix}\]

\pause
\vspace{0.3em}

范数最小的最小二乘解:
\[\mathbf{x}^+ = A^+\mathbf{b} = \begin{pmatrix}\frac{1}{4}&\frac{1}{4}&0\\\frac{1}{4}&\frac{1}{4}&0\\0&0&1\end{pmatrix}\begin{pmatrix}0\\4\\2\end{pmatrix} = \begin{pmatrix}1\\1\\2\end{pmatrix}, \qquad \lVert\mathbf{x}^+\rVert^2 = 6.\]

\pause
\vspace{0.3em}

\textbf{对比}: $(0,2,2)^T$ 和 $(2,0,2)^T$ 均为最小二乘解, 但 $\lVert(0,2,2)^T\rVert^2 = \lVert(2,0,2)^T\rVert^2 = 8 > 6$.

\vspace{0.3em}
\textbf{几何意义}: $\mathbf{x}^+ = (1,1,2)^T$ 是仿射子空间 $\{x_1+x_2=2,\, x_3=2\}$ 上距原点最近的点; 它在零空间方向 $(1,-1,0)^T$ 上分量为零, 完全落在行空间 $C(A^T)$ 中.

\end{frame}

%----------------------------------------------------------------------------------------

\end{document}
